\documentclass[10pt,twocolumn]{article}
\usepackage[margin=1in]{geometry}
\usepackage{times}
\usepackage{graphicx}
\usepackage{amsmath}
\usepackage{amssymb}
\usepackage{booktabs}
\usepackage{hyperref}
\usepackage{xcolor}
\usepackage{microtype}
\usepackage{enumitem}

\hypersetup{
    colorlinks=true,
    linkcolor=blue,
    citecolor=blue,
    urlcolor=blue
}

\title{\textbf{SynID: Zero-Shot Identity-Consistent Image Generation\\
via Synthetic Bootstrapping and On-the-Fly UNet Adaptation}}

\author{
Anonymous Author(s)\\
\textit{Under review}
}

\date{}

\begin{document}

\maketitle

% ─────────────────────────────────────────────────────────────
\begin{abstract}
We present \textbf{SynID}, a novel framework for identity-consistent image generation that requires no real reference images, no external dataset, and no large-scale pretraining. Given only a text description of a character, SynID generates multiple images of the same identity across diverse expressions, scenes, and poses. The system operates through a closed-loop self-distillation pipeline: synthetic anchor images are generated from the model's own prior, a multi-token identity projector is trained on the resulting CLIP embeddings, bootstrap refinement selects high-quality candidates to refine the identity embedding, and a lightweight UNet cross-attention adapter is trained on-the-fly using only eight synthetic images in approximately 260 steps. Identity is injected at two levels simultaneously---text embedding space for coarse alignment and UNet cross-attention for fine-grained facial consistency. On five diverse character types, SynID achieves a mean face-weighted CLIP identity score of \textbf{0.954} and an ArcFace score of \textbf{0.791}, comparable to Arc2Face ($\sim$0.79 on FFHQ) while requiring zero real images. In a controlled comparison, SynID outperforms IP-Adapter FaceID on CLIP identity (0.969 vs.\ 0.854). The entire pipeline runs in approximately five minutes per character on a single T4 GPU. The released implementation is packaged as an app-first backend with staged model loading and an interactive interface layered on top of the research pipeline.
\end{abstract}

% ─────────────────────────────────────────────────────────────
\section{Introduction}

Identity-consistent image generation---producing multiple images of the same person across different contexts---is a fundamental challenge in generative modeling. Existing approaches fall into two categories: fine-tuning methods (DreamBooth~\cite{ruiz2023dreambooth}, LoRA~\cite{hu2022lora}) that require multiple real photographs and significant compute, and adapter-based methods (IP-Adapter~\cite{ye2023ipadapter}, InstantID~\cite{wang2024instantid}) that require a real reference image at inference time and are trained on millions of image pairs.

Both paradigms share a critical dependency: \textit{real images of the target identity}. This limits their applicability in scenarios where no reference image exists---character design, synthetic data generation, privacy-preserving personalization, and creative applications where the user wants to generate a character from a text description alone.

We ask a different question: \textit{can identity-consistent generation be achieved from text alone, with no real images at any stage?}

SynID answers yes. The key insight is that modern diffusion models already contain rich identity priors---they can generate consistent-looking characters from text descriptions. SynID exploits this prior through a closed-loop self-distillation process: the model generates its own training data, refines its own identity embedding, and trains its own adapter, all without any external supervision. We call this paradigm \textbf{bootstrapped synthetic identity}.

Our contributions are:
\begin{enumerate}[leftmargin=*,topsep=2pt,itemsep=1pt]
    \item \textbf{Synthetic identity formation}: Identity is initialized from the model's own generative prior via a softmax-weighted multi-anchor ensemble, producing a robust CLIP embedding without any real images. No prior work generates fresh synthetic identities purely from text and averages CLIP image embeddings in this way.
    \item \textbf{Closed-loop self-distillation}: A bootstrap refinement loop generates expression-diverse candidates, scores them by a balanced CLIP+ArcFace metric with diversity enforcement, and iteratively refines the identity embedding. No prior identity personalization method has explicitly generated new samples, scored them by CLIP+ArcFace, and re-blended with anchors to retrain.
    \item \textbf{On-the-fly UNet adaptation}: A lightweight cross-attention adapter is trained in $\sim$260 steps on eight synthetic images, injecting identity directly into the UNet's denoising process. Training adapters on synthetic images per-identity on-the-fly, without any offline pretraining, is novel.
    \item \textbf{Dual-level identity injection}: Identity tokens are injected at both the text embedding level (coarse) and UNet cross-attention level (fine-grained). Prior methods use only one level---IP-Adapter uses only cross-attention, textual inversion uses only text.
    \item \textbf{Self-correcting drift correction}: A feedback loop measures CLIP drift between generated images and the identity embedding, fine-tuning the projector to close the gap. No published identity method has used a probe-image-to-measure-embedding-shift mechanism of this kind.
    \item \textbf{Identity-aware negative conditioning}: The negative prompt embedding is explicitly pushed away from the identity direction, actively suppressing identity drift during classifier-free guidance.
\end{enumerate}

% ─────────────────────────────────────────────────────────────
\section{Related Work}

\paragraph{Fine-tuning approaches.}
DreamBooth~\cite{ruiz2023dreambooth} and LoRA~\cite{hu2022lora} fine-tune diffusion model weights on 3--20 real images of a subject. These methods achieve high identity fidelity but require real photographs and significant compute (15--30 minutes per subject). ID-Booth~\cite{tomasevic2025idbooth} extends this with a triplet identity loss for improved consistency, but still requires real images.

\paragraph{Adapter-based approaches.}
IP-Adapter~\cite{ye2023ipadapter} trains a cross-attention adapter on $\sim$1M image-text pairs, enabling image-prompted generation at inference time. IP-Adapter FaceID extends this with ArcFace embeddings for face-specific conditioning. InstantID~\cite{wang2024instantid} combines face embeddings with ControlNet for high-fidelity identity transfer. Inv-Adapter~\cite{hong2024invadapter} inverts a single reference image and trains a per-identity adapter. All these methods require at least one real reference image at inference time. F-Bench~\cite{liu2025fbench} provides a large human-evaluated benchmark comparing these methods, finding FaceID-Plus achieves the best overall performance.

\paragraph{Stacked ID embedding.}
PhotoMaker~\cite{li2024photomaker} personalizes generation by encoding multiple real ID images into a stacked ID embedding, achieving ArcFace similarity of $\sim$0.618 on their evaluation set.

\paragraph{Foundation models for faces.}
Arc2Face~\cite{papantoniou2024arc2face} trains a face foundation model on 21M facial images, achieving ArcFace similarity of $\sim$0.79 on a held-out FFHQ identity set. It requires a real ArcFace embedding as input.

\paragraph{Concurrent zero-shot personalization.}
EZIGen~\cite{zhang2024ezigen} achieves zero-shot personalized image generation but requires a real subject image as input. Inv-Adapter~\cite{hong2024invadapter} similarly requires one real image for inversion.

\paragraph{Key difference.}
SynID requires no real images at any stage---not for training, not for inference. The identity is entirely synthetic, initialized from the diffusion model's own prior and refined through self-distillation. This distinguishes it from all methods above. The closest conceptual relatives are identity token methods~\cite{chen2023identitytokens} and DreamBooth-style iterative refinement, but neither operates in a fully synthetic, closed-loop manner.

% ─────────────────────────────────────────────────────────────
\section{Method}

\subsection{Overview}

Given a text prompt $p$ describing a character, SynID produces a \texttt{CharacterProfile} containing identity tokens $\mathbf{z} \in \mathbb{R}^{4 \times 768}$ that can be used to generate arbitrarily many consistent images. The pipeline consists of five stages: (1) multi-anchor ensemble embedding, (2) multi-token projector training, (3) bootstrap refinement, (4) drift correction, and (5) on-the-fly UNet adapter training.

\subsection{Multi-Anchor Ensemble Embedding}

A single synthetic anchor image may be unrepresentative due to seed-dependent variation. We generate $N=4$ anchor images from the same prompt with different seeds and compute a softmax-weighted ensemble embedding:

\begin{equation}
\mathbf{e}_{ens} = \text{normalize}\!\left(\sum_{i=1}^{N} w_i \cdot \text{CLIP}(I_i^{anc})\right)
\end{equation}

where $w_i = \text{softmax}_i\!\left(\cos\!\left(\text{CLIP}(I_i^{anc}),\, \text{CLIP}_{text}(p)\right)\right)$. Anchors more similar to the text prompt receive higher weight.

We use a face-weighted CLIP encoding that blends full-image features (40\%) with face-crop features (60\%):

\begin{equation}
\text{CLIP}_{fw}(I) = \text{normalize}\!\left(0.4\,\text{CLIP}_{full}(I) + 0.6\,\text{CLIP}_{crop}(I)\right)
\end{equation}

where $\text{CLIP}_{crop}$ encodes the center face region $(x \in [0.18w, 0.82w],\; y \in [0.03h, 0.67h])$, targeting the facial region without requiring a face detector.

\subsection{Multi-Token Identity Projector}

We train a lightweight MLP projector $f_\theta: \mathbb{R}^{512} \rightarrow \mathbb{R}^{4 \times 768}$ (architecture: Linear(512$\to$1024) $\to$ SiLU $\to$ Linear(1024$\to$4$\times$768)) that maps the CLIP image embedding to four identity tokens in the SD text encoder's embedding space. The training objective is:

\begin{equation}
\mathcal{L}_{proj} = \mathcal{L}_{text} + \mathcal{L}_{pt} + \lambda_{div}\mathcal{L}_{div} + \lambda_{norm}\mathcal{L}_{norm}
\end{equation}

where $\mathcal{L}_{text} = 1 - \cos(\bar{\mathbf{z}}, \mathbf{t})$ aligns the mean token $\bar{\mathbf{z}}$ to the masked-mean text embedding $\mathbf{t}$, and $\mathcal{L}_{pt} = 0.3 \cdot \frac{1}{K}\sum_k (1 - \cos(\mathbf{z}_k, \mathbf{t}))$ aligns each token individually. $\mathcal{L}_{div}$ encourages token diversity and $\mathcal{L}_{norm}$ regularizes norms toward 1.0. We use $\lambda_{div}=0.05$, $\lambda_{norm}=0.01$, Adam optimizer, 250 steps.

When ArcFace embeddings are available, an additional face verification loss is added:
\begin{equation}
\mathcal{L}_{arc} = 0.3 \cdot \left(1 - \cos\!\left(\bar{\mathbf{z}},\, \mathbf{W}_{proj}\,\mathbf{e}_{arc}\right)\right)
\end{equation}
where $\mathbf{W}_{proj} \in \mathbb{R}^{768 \times 512}$ projects the ArcFace embedding to text space.

\subsection{Bootstrap Refinement}

The initial projector generates $M=8$ bootstrap images with diverse expression prompts. Each candidate is scored:

\begin{equation}
s_i = 0.45\,\text{CLIP}_{id}(I_i) + 0.35\,\text{ArcFace}(I_i) + 0.20\,\text{CLIP}_{prompt}(I_i)
\end{equation}

Candidates are selected with diversity enforcement: a candidate is included only if its identity similarity exceeds 0.85 and its pairwise similarity to already-selected candidates is below 0.92. The top-$K=2$ candidates are blended with the anchor:

\begin{equation}
\mathbf{e}_{ref} = \text{normalize}\!\left(0.55\,\mathbf{e}_{ens} + 0.45\,\overline{\mathbf{e}_{sel}}\right)
\end{equation}

The projector is retrained on $\mathbf{e}_{ref}$ for 150 steps.

\subsection{Drift Correction}

A self-correcting feedback loop closes the gap between the projector's target and the system's actual output:

\begin{enumerate}[leftmargin=*,topsep=2pt,itemsep=1pt]
    \item Generate one image $I_{gen}$ using current identity tokens
    \item Compute drift: $d = 1 - \cos(\text{CLIP}_{fw}(I_{gen}),\, \mathbf{e}_{ref})$
    \item If $d > 0.015$: blend $\mathbf{e}_{ref} \leftarrow \text{normalize}(0.65\,\mathbf{e}_{ref} + 0.35\,\text{CLIP}_{fw}(I_{gen}))$
    \item Fine-tune projector 40 steps on blended embedding
    \item Repeat for 2 rounds
\end{enumerate}

\subsection{UNet Cross-Attention Adapter}

We attach a lightweight cross-attention adapter to every transformer block in the UNet's down, mid, and up blocks. The adapter uses identity tokens as key-value pairs and UNet hidden states as queries:

\begin{equation}
\text{Adapter}(\mathbf{h}, \mathbf{z}) = \mathbf{h} + \alpha \cdot \text{CrossAttn}_Q(\mathbf{h})\,\text{CrossAttn}_{KV}(\mathbf{z})
\end{equation}

We use a graduated per-block scale: down blocks ($\alpha = 0.25$), mid block ($\alpha = 0.375$), up blocks ($\alpha = 0.5$). The output projection is zero-initialized, ensuring the adapter starts as a no-op.

\paragraph{On-the-fly training.}
The adapter is trained on the 8 bootstrap images with a three-term objective:

\begin{equation}
\mathcal{L}_{adp} = \mathcal{L}_{MSE} + \lambda_{c}\mathcal{L}_{CLIP} + \lambda_{a}\mathcal{L}_{ArcFace}
\end{equation}

$\mathcal{L}_{MSE}$ is the standard diffusion noise prediction loss. $\mathcal{L}_{CLIP}$ and $\mathcal{L}_{ArcFace}$ are computed on decoded predictions every 10 and 25 steps respectively. Scale warmup ramps $\alpha$ from 0.1 to target over the first 50\% of training. We use AdamW with cosine LR schedule, 200 steps ($\lambda_c=0.1$), followed by a 60-step ArcFace fine-tune pass ($\lambda_a=0.2$).

\subsection{Identity Injection}

At generation time, identity tokens are injected at two levels:

\paragraph{Text embedding level (coarse).}
\begin{equation}
\mathbf{c}[:, -N:, :] \mathrel{+}= \alpha_{adp} \cdot \mathbf{z}
\end{equation}
where $\alpha_{adp}$ scales adaptively with prompt complexity.

\paragraph{UNet level (fine-grained).}
The adapter injects identity into every transformer block via cross-attention, influencing the entire denoising trajectory.

\paragraph{Identity-aware negative conditioning.}
\begin{equation}
\mathbf{c}_{neg}[:, -N:, :] \mathrel{-}= 0.5 \cdot \mathbf{z}
\end{equation}
This actively suppresses identity drift rather than passively zeroing negative slots.

\subsection{Reference Implementation and Interface}

The released system separates \textbf{backend initialization} from \textbf{user interaction}. Model definitions load immediately, but heavyweight assets---DreamShaper, ControlNet, CLIP, and OpenPose---are initialized only through an explicit backend setup step. This keeps startup responsive and makes the system usable both as a library and as an application. On top of the backend, we provide an interactive UI that exposes staged pipeline loading, character creation, generation, evaluation, checkpointing, and mobile access through a QR linked to the public app URL. These packaging choices are not the research contribution themselves, but they make the full pipeline reproducible and usable beyond a notebook-style workflow.

% ─────────────────────────────────────────────────────────────
\section{Experiments}

\subsection{Setup}

\textbf{Base model:} Lykon/DreamShaper (SD 1.5).
\textbf{CLIP:} openai/clip-vit-large-patch14 when available, with fallback to openai/clip-vit-base-patch32.
\textbf{Face verification:} InsightFace buffalo\_sc.
\textbf{Pose:} ControlNet OpenPose.
\textbf{Hardware:} NVIDIA T4 GPU.
\textbf{Time per character:} $\sim$5 minutes.

\subsection{Evaluation Metrics}

\begin{itemize}[leftmargin=*,topsep=2pt,itemsep=1pt]
    \item \textbf{Face-weighted CLIP (FW-CLIP):} Blend of full-image (40\%) and face-crop (60\%) CLIP cosine similarity to anchor embedding.
    \item \textbf{ArcFace:} InsightFace buffalo\_sc cosine similarity between generated and anchor face embeddings.
    \item \textbf{Pairwise consistency:} Mean pairwise FW-CLIP similarity across 8 generated images.
\end{itemize}

\subsection{Multi-Character Results}

\begin{table}[h]
\centering
\caption{Identity consistency across five diverse character types.}
\label{tab:multichar}
\begin{tabular}{lcc}
\toprule
Character & FW-CLIP & Pairwise \\
\midrule
Woman (brunette) & 0.9515 & 0.9435 $\pm$ 0.022 \\
Elderly man & 0.9508 & 0.9430 $\pm$ 0.031 \\
Anime girl & 0.9655 & 0.9541 $\pm$ 0.013 \\
Young man & 0.9407 & 0.9528 $\pm$ 0.019 \\
Woman (redhead) & 0.9625 & 0.9411 $\pm$ 0.013 \\
\midrule
\textbf{Mean} & \textbf{0.9542} & \textbf{0.9469 $\pm$ 0.020} \\
\bottomrule
\end{tabular}
\end{table}

\subsection{Evaluation Suites}

\begin{table}[h]
\centering
\caption{Evaluation across four axes (woman redhead character).}
\label{tab:suites}
\begin{tabular}{lcccc}
\toprule
Suite & CLIP & $\sigma$ & ArcFace & $\sigma$ \\
\midrule
Expression & 0.9523 & 0.021 & 0.5113 & 0.071 \\
Scene & 0.9349 & 0.013 & 0.5502 & 0.022 \\
Pose & 0.9635 & 0.011 & 0.5882 & 0.050 \\
Seed & 0.9628 & 0.005 & 0.6105 & 0.044 \\
\bottomrule
\end{tabular}
\end{table}

The seed suite (std=0.005) demonstrates strong robustness across random seeds. The pose suite uses reduced ControlNet scale (0.35) to allow natural pose variation.

\subsection{Ablation Study}

\begin{table}[h]
\centering
\caption{Ablation: contribution of each component.}
\label{tab:ablation}
\begin{tabular}{lccc}
\toprule
Configuration & CLIP & ArcFace & Time (s) \\
\midrule
Baseline & 0.9603 & 0.7562 & 55 \\
+ Ensemble & 0.9676 & 0.7581 & 78 \\
+ Bootstrap + Drift & 0.9665 & 0.7560 & 154 \\
\textbf{+ UNet Adapter} & \textbf{0.9690} & \textbf{0.7912} & 158 \\
\bottomrule
\end{tabular}
\end{table}

The UNet adapter provides the largest ArcFace improvement (+3.5\%), confirming that cross-attention injection captures facial geometry that text-space injection alone cannot.

\subsection{Comparison with Prior Work}

\begin{table}[h]
\centering
\caption{Comparison with identity-conditioning methods. Benchmarks are not standardized across methods; direct numeric comparisons should be interpreted with this caveat.}
\label{tab:comparison}
\begin{tabular}{lcccc}
\toprule
Method & CLIP & ArcFace & Real img & Data \\
\midrule
PhotoMaker~\cite{li2024photomaker} & --- & $\sim$0.618$^\ddagger$ & Yes & Real \\
IP-Adapter FaceID & 0.854 & 0.132$^\dagger$ & Yes & $\sim$1M \\
Arc2Face~\cite{papantoniou2024arc2face} & --- & $\sim$0.79$^\ddagger$ & Yes & 21M \\
\textbf{SynID (ours)} & \textbf{0.969} & \textbf{0.791} & \textbf{No} & \textbf{8 syn.} \\
\bottomrule
\end{tabular}
\end{table}

$^\dagger$IP-Adapter run without pose conditioning; scores may not be directly comparable. $^\ddagger$Reported on authors' own evaluation sets.

SynID achieves ArcFace scores comparable to Arc2Face---a model trained on 21M real faces---while requiring zero real images and no pretraining.

% ─────────────────────────────────────────────────────────────
\section{Discussion}

\paragraph{What the system learns.}
SynID does not learn identity in the same sense as face recognition systems. Rather, it learns to steer the diffusion model's existing identity priors toward a specific point in the generative manifold. The projector maps a CLIP embedding to a text-space direction that activates the model's internal representation of the described character. The adapter reinforces this direction at every denoising step.

\paragraph{Limitations.}
(1) \textit{Pose dependence}: all generated images share the same head angle from the anchor pose skeleton. (2) \textit{Expression sensitivity}: ArcFace scores drop from 0.61 (seed) to 0.51 (expression), as strong expressions change facial geometry. (3) \textit{Base model dependence}: identity consistency is bounded by the base model's generative prior. (4) \textit{Non-face subjects}: the system is designed for portrait generation and does not generalize to animals or full-body characters.

\paragraph{Failure cases.}
Extreme pose changes (profile view, $>$45° tilt), very similar character descriptions, and non-face subjects cause identity drift or failure.

\paragraph{Ethical considerations.}
SynID generates entirely synthetic identities from text descriptions and does not reproduce real people's likenesses. However, the system could be used to generate persistent synthetic personas at scale, raising concerns about synthetic identity misuse in disinformation or fraud contexts. Responsible deployment should include watermarking or provenance tracking for generated images.

% ─────────────────────────────────────────────────────────────
\section{Conclusion}

We presented SynID, a zero-shot identity-consistent generation system requiring no real reference images at any stage. Through a closed-loop self-distillation pipeline combining multi-anchor ensemble embedding, bootstrap refinement, drift correction, and on-the-fly UNet adaptation, SynID achieves identity consistency comparable to systems trained on millions of real images, while operating entirely within the model's own generative prior. The core finding is that \textbf{bootstrapped synthetic identity}---initializing from the model's own prior and refining through self-distillation---can replace dataset-driven identity conditioning for text-described characters. Future work includes extending to SDXL, removing ControlNet pose dependence, and multi-character consistency.

% ─────────────────────────────────────────────────────────────
\bibliographystyle{plain}
\begin{thebibliography}{99}

\bibitem{ruiz2023dreambooth}
N.~Ruiz, Y.~Li, V.~Jampani, Y.~Pritch, M.~Rubinstein, and K.~Aberman.
\newblock {DreamBooth}: Fine tuning text-to-image diffusion models for subject-driven generation.
\newblock In \textit{CVPR}, 2023.

\bibitem{hu2022lora}
E.~J. Hu, Y.~Shen, P.~Wallis, Z.~Allen-Zhu, Y.~Li, S.~Wang, L.~Wang, and W.~Chen.
\newblock {LoRA}: Low-rank adaptation of large language models.
\newblock In \textit{ICLR}, 2022.

\bibitem{ye2023ipadapter}
H.~Ye, J.~Zhang, S.~Liu, X.~Han, and W.~Yang.
\newblock {IP-Adapter}: Text compatible image prompt adapter for text-to-image diffusion models.
\newblock \textit{arXiv:2308.06721}, 2023.

\bibitem{wang2024instantid}
Q.~Wang, X.~Bai, H.~Wang, Z.~Qin, A.~Chen, and W.~Sun.
\newblock {InstantID}: Zero-shot identity-preserving generation in seconds.
\newblock \textit{arXiv:2401.07519}, 2024.

\bibitem{papantoniou2024arc2face}
F.~Papantoniou, A.~Lattas, S.~Moschoglou, J.~Deng, B.~Kainz, and S.~Zafeiriou.
\newblock {Arc2Face}: A foundation model for {ID}-consistent human faces.
\newblock In \textit{ECCV}, 2024.

\bibitem{deng2019arcface}
J.~Deng, J.~Guo, N.~Xue, and S.~Zafeiriou.
\newblock {ArcFace}: Additive angular margin loss for deep face recognition.
\newblock In \textit{CVPR}, 2019.

\bibitem{zhang2023controlnet}
L.~Zhang, A.~Rao, and M.~Agrawala.
\newblock Adding conditional control to text-to-image diffusion models.
\newblock In \textit{ICCV}, 2023.

\bibitem{rombach2022ldm}
R.~Rombach, A.~Blattmann, D.~Lorenz, P.~Esser, and B.~Ommer.
\newblock High-resolution image synthesis with latent diffusion models.
\newblock In \textit{CVPR}, 2022.

\bibitem{li2024photomaker}
Z.~Li, M.~Cao, X.~Wang, Z.~Qi, M.-M. Cheng, and Y.~Shan.
\newblock {PhotoMaker}: Customizing realistic human photos via stacked {ID} embedding.
\newblock In \textit{CVPR}, 2024.

\bibitem{tomasevic2025idbooth}
D.~Tomašević, M.~Donoser, and L.~Bossard.
\newblock {ID-Booth}: Identity-consistent face generation with diffusion models.
\newblock \textit{arXiv:2504.07392}, 2025.

\bibitem{hong2024invadapter}
Y.~Hong, M.~Zhang, J.~Gu, M.~Fu, L.~Wang, L.~Wei, and H.~Zhu.
\newblock {Inv-Adapter}: {ID} customization generation via image inversion.
\newblock \textit{arXiv:2406.02881}, 2024.

\bibitem{liu2025fbench}
Y.~Liu et al.
\newblock {F-Bench}: Rethinking human preference evaluation metrics for benchmarking face generation, customization, and restoration.
\newblock \textit{arXiv:2412.13155}, 2025.

\bibitem{zhang2024ezigen}
Z.~Zhang et al.
\newblock {EZIGen}: Enhancing zero-shot subject-driven image generation with precise subject encoding and decoupled guidance.
\newblock \textit{arXiv}, 2024.

\bibitem{chen2023identitytokens}
X.~Chen et al.
\newblock Identity-preserving text-to-image generation using identity tokens.
\newblock \textit{arXiv}, 2023.

\end{thebibliography}

\end{document}
